% Generated draft: requires mathematical and source review.
\documentclass[11pt]{article}
\usepackage[margin=25mm]{geometry}
\usepackage{fontspec}
\setmainfont{Latin Modern Roman}
\usepackage{amsmath,amssymb,mathtools,unicode-math}
\setmathfont{Latin Modern Math}
\usepackage{enumitem,needspace,xurl,hyperref}
\setlength{\parindent}{0pt}
\setlength{\parskip}{5pt}
\setlength{\emergencystretch}{3em}
\newcommand{\N}{\mathbb{N}}
\newcommand{\Z}{\mathbb{Z}}
\newcommand{\Q}{\mathbb{Q}}
\newcommand{\R}{\mathbb{R}}
\newcommand{\C}{\mathbb{C}}
\newcommand{\F}{\mathbb{F}}
\newcommand{\A}{\mathbb{A}}
\newcommand{\PP}{\mathbb{P}}
\newcommand{\E}{\mathbb{E}}
\newcommand{\eps}{\varepsilon}
\newcommand{\dd}{\,\mathrm{d}}
\begin{document}
\section*{UnsolvedMath ID: 30006450}
\textbf{Source problem code:} catalogue entry
\subsection{Definitions and mathematical statement}
Let \(G=(V,E)\) be a graph. To obtain the finite generating polynomial displayed in the source, suppose that \(E\) is finite; this finiteness hypothesis is implicit rather than explicit in the supplied record. Define
\[
\mathcal F(G)=\{F\subseteq E:(V,F)\text{ is acyclic}\},\qquad
x^F=\prod_{e\in F}x_e,\qquad
p_{\mathcal F}(x)=\sum_{F\in\mathcal F(G)}x^F,
\]
where the empty product is \(1\). Thus \(p_{\mathcal F}\) is a multiaffine polynomial with nonnegative coefficients. Here every forest is regarded as a subgraph on the full vertex set, allowing isolated vertices; no maximality or connectivity requirement is imposed. This follows the original excerpt's description of forests as subsets of \(E\). The graph conventions concerning loops and parallel edges are not specified.

\textbf{Conjecture 6 (attributed in the supplied excerpt to Kahn and Grimmett--Winkler).} For every graph in the intended domain, the source asks whether
\[
R(1,\{1\})\text{ holds for }p_{\mathcal F}.
\]
The supplied record calls this a negative-dependence property but does not define \(R(c,S)\). Consequently, its inequalities, quantified indices, evaluation points, and any conditioning or closure requirements cannot be recovered from the supplied evidence. Even whether the symbol \(1\) denotes the all-ones evaluation vector is not explicitly defined. The displayed conjecture therefore preserves the source notation and remains incomplete as a self-contained mathematical assertion.

\textbf{Broader proposed approach.} The original excerpt proposes finding a polynomial class \(\mathcal C\), possibly containing the real-stable polynomials and contained in the Lorentzian polynomials, or in a closely related class, such that, for some desired constant \(c\in[1,2]\), every polynomial in \(\mathcal C\) satisfies the source property
\[
R(c,\R_{\geq0}^{\,n}),
\]
with \(n\) the relevant number of variables. The stronger wording in the cleaned statement asks for a class between the real-stable and Lorentzian classes; the original excerpt presents these containments as possibilities. The ambient polynomial domain, degree and homogeneity conventions, meaning of a closely related class, and definition of \(R\) are not supplied, so no precise class-existence assertion can be fixed. The excerpt also refers to two open problems without supplying the other problem.
\subsection{Short English statement}
The conjecture concerns the polynomial enumerating all acyclic edge subsets of a graph, including the empty forest and forests with isolated vertices. It asks whether this polynomial has the source's negative-dependence property \(R(1,\{1\})\). A broader suggestion seeks a polynomial class connecting real stability and Lorentzian behavior with a constant \(c\in[1,2]\). The missing definition of \(R\) and unspecified class conventions prevent a fully self-contained formulation.
\subsection{Sources}
\begin{itemize}
\item Supplied local UnsolvedMath registry record, permanent ID \(30006450\), title: \emph{Negative Dependence for Random Spanning Forests}; exact source problem code: \texttt{catalogue entry}.
\item Supplied report citation: \emph{Arbeitsgemeinschaft: Combinatorial Hodge Theory (2026)}, Oberwolfach Reports, pp.~2545--2592, DOI \texttt{10.4171/owr/2025/47}. Source URL and supplied literature-source link, labeled Oberwolfach Report 47/2025: \url{https://doi.org/10.4171/owr/2025/47}. This citation is supplied evidence; it was not independently accessed or verified for this draft.
\item The original excerpt labels the forest assertion Conjecture~6 and attributes it to Kahn and Grimmett--Winkler, with a reference to [4]. The bibliographic entry for [4] is absent from the supplied record.
\item Literature assessment dated exactly \texttt{2026-08-22}: supplied dated curation labels the problem open and states that the report presents the proposed interpolating class as a general approach, supplying no proof of its existence. The background describes this as dated literature triage derived from supplied OWR curation and calls for independent verification of time-sensitive status and cited claims. This assessment is not independently verified here.
\item The registry labels the statement \texttt{corrected\_verified} and reports that the official report was checked and its notation repaired. These are supplied provenance claims, not verification performed for this draft. The record gives proposed year 2025, while the report title includes 2026 and the report identifier includes 2025; their relationship is not explained.
\end{itemize}
\end{document}
